\chapter{Review of Literature}

\section{Introduction}

The review of literature provides the conceptual foundation for the present study. It helps position the research within existing academic and managerial discussions related to sentiment analysis, online customer reviews, app-based service platforms, Natural Language Processing, and brand perception. Since the present study examines Zomato customer sentiment through Google Play Store reviews, the literature reviewed in this chapter focuses on how digital text data can be used to understand customer experience, service quality issues, and brand-related insights \cite{pang2008,liu2012,lemon2016}.

\section{Sentiment Analysis and Customer Opinion Mining}

Sentiment analysis is a widely used text analytics approach for identifying the emotional orientation expressed in textual data. In business and marketing contexts, it is commonly used to classify customer opinions into positive, negative, and neutral categories. The growing use of sentiment analysis reflects the increasing importance of digital feedback as a source of market intelligence. Organizations use sentiment analysis to understand how customers react to products, services, and brand interactions across online platforms \cite{pang2008,liu2012}.

Opinion mining is closely related to sentiment analysis and focuses on extracting evaluative information from textual data. In practice, sentiment analysis enables firms to process large volumes of customer comments more efficiently than manual review. This is especially valuable when customer feedback is continuous, unstructured, and generated at scale. For service-oriented platforms, sentiment analysis supports the identification of recurring strengths and weaknesses in the customer journey \cite{pang2008,liu2012,hutto2014}.

\section{Online Customer Reviews as a Source of Business Insight}

Online customer reviews are an important form of electronic word-of-mouth because they influence public perception and shape purchase decisions. Customers rely on reviews to evaluate service quality, trustworthiness, value for money, and overall satisfaction before making platform choices. At the same time, firms can use reviews to understand consumer expectations and identify areas requiring service improvement \cite{chevalier2006,mudambi2010,filieri2015,park2009}.

For digital platforms, review data offers an especially rich source of insight because users often describe specific service encounters in detail. Unlike aggregate ratings alone, textual reviews reveal the reasons behind satisfaction or dissatisfaction. They help firms distinguish between issues such as delayed delivery, poor customer support, technical problems in the app, pricing concerns, and refund-related dissatisfaction. As a result, online reviews serve both as a feedback mechanism and as a strategic input for managerial decision-making \cite{mudambi2010,filieri2015}.

\section{App Reviews in Platform-Based Service Businesses}

Application reviews posted on digital app stores differ from generic product reviews because they often combine comments on both service performance and technological usability. In the case of food delivery and restaurant aggregation platforms, users evaluate not only the quality of the service outcome but also the usability of the mobile application through which the service is accessed. This makes app review data especially valuable for businesses such as Zomato, where customer experience is shaped by both operational execution and digital interface quality \cite{pagano2013,guzman2017,martin2017}.

Google Play Store reviews are particularly useful for analysis because they include both star ratings and written feedback. This dual structure allows researchers to interpret ratings as a broad measure of satisfaction while using review text to identify the exact themes driving customer sentiment. For the present study, this feature is important because the review rating serves as a sentiment proxy, while the review text provides richer thematic evidence regarding customer concerns and expectations \cite{pagano2013,martin2017}.

\section{Natural Language Processing in Review Analytics}

Natural Language Processing makes it possible to convert large volumes of human language into structured and analyzable information. In review analytics, NLP techniques are used for cleaning text, reducing noise, standardizing language, classifying sentiment, and identifying recurring issues or themes. This is particularly useful in digital environments where data is generated continuously and where manual interpretation becomes impractical \cite{manning1999,jurafsky2024}.

NLP-based analysis is valuable not only because it supports automation, but also because it improves consistency in how large datasets are interpreted. In practical applications, NLP techniques are often combined with rating information, keyword rules, TF-IDF style vector representations, or machine learning methods depending on the nature of the dataset and the purpose of the study \cite{salton1988,hosmer2013}. In the present project, NLP is relevant because it supports preprocessing, sentiment categorization, and theme identification in a large corpus of Google Play Store reviews related to Zomato.

\section{Themes Influencing Customer Perception in Food Delivery Apps}

Customer perception of food delivery platforms is influenced by a combination of service quality, reliability, responsiveness, convenience, and trust. In digital food-service platforms, customers commonly evaluate the timeliness of delivery, cancellation and refund handling, app performance, pricing and fee transparency, customer support responsiveness, and the overall ease of placing and tracking orders. These operational and experiential factors strongly influence satisfaction and dissatisfaction in the broader customer journey \cite{lemon2016,verhoef2010}.

Thematic analysis of customer reviews helps identify which of these issues dominate public perception. In the context of app-based services, theme-level analysis is useful because overall sentiment alone does not fully explain the sources of customer satisfaction or frustration. A platform may receive positive sentiment on one dimension, such as app usability, while also attracting negative sentiment on another, such as refund processing. Therefore, theme identification complements sentiment analysis by providing more actionable managerial insight \cite{filieri2015,lemon2016}.

\section{Digital Reviews, Brand Strategy, and Customer Experience}

Brand perception is increasingly shaped by digital interactions and publicly visible customer feedback. Reviews posted on app stores and other platforms do not merely reflect private customer opinion; they also influence how potential users perceive the brand before trying the service themselves. This means that customer reviews have both diagnostic value and reputational value. A concentration of negative reviews around certain themes can affect trust, trial, loyalty, and recommendation behavior \cite{chevalier2006,keller2013}.

From a managerial perspective, customer review analytics supports both brand strategy and customer experience management. When firms identify the themes most strongly associated with negative sentiment, they can prioritize corrective action in areas that most affect customer perception. Similarly, themes associated with positive feedback help reveal the strengths that should be reinforced in communication and service positioning. For Zomato, such insights are valuable because they can connect user feedback with decisions related to service quality, issue resolution, app design, and competitive positioning \cite{lemon2016,verhoef2010,keller2013}.

\section{Prior Zomato-Focused Benchmark Study}

A directly relevant prior study for the present research is the paper titled \textit{Sentiment Analysis of Customer Reviews in Zomato Bangalore Restaurants Using Random Forest Classifier} by Jonathan, Sihotang, and Martin \cite{jonathan2019}. That study is useful because it applies sentiment classification specifically to Zomato-related review data and uses Random Forest as the core modeling approach. It therefore provides a practical benchmark for comparing methodological direction and model performance.

At the same time, the prior paper is narrower in scope than the present study. Its emphasis is primarily on classification performance, whereas the current project extends the discussion toward theme-level issue analysis, managerial interpretation, dashboard visualization, and competitor comparison. The benchmark paper is therefore valuable not only as a literature source, but also as a reference point for explaining how the current study broadens a model-focused approach into a business decision-support approach.

\section{Research Framework of the Study}

Based on the literature reviewed, the present study adopts a framework in which customer reviews act as the primary source of digital feedback, sentiment analysis captures the overall emotional orientation of customer opinion, and thematic classification identifies the major issues shaping perception. These analytical outputs are then interpreted in terms of their implications for customer experience and brand strategy.

The framework of the study can be summarized as follows:

\begin{enumerate}[label=\arabic*.]
    \item Google Play Store reviews provide the raw customer feedback data.
    \item Review ratings indicate broad customer sentiment as positive, negative, or neutral.
    \item Review text reveals the dominant themes affecting customer perception.
    \item Theme-wise sentiment patterns help identify major strengths and problem areas.
    \item The findings support managerial recommendations related to brand strategy and customer experience improvement.
\end{enumerate}

\section{Chapter Summary}

This chapter reviewed the literature relevant to sentiment analysis, online reviews, app-based service businesses, Natural Language Processing, customer perception, and brand strategy. The review establishes the conceptual basis for analyzing Zomato's Google Play Store reviews as a source of customer insight. The next chapter presents the research methodology adopted for data collection, preprocessing, labeling logic, and analytical design.
